\begin{abstract}
For an elementary rule $f$, the second-order automaton $F_f$ on $(\{0,1\}^2)^{\Z}$, given by $y_{t+1}=y_{t-1}\oplus f(y_t)$, is reversible, and the automorphism group of its spacetime subshift is the centralizer of $F_f$ in the automorphism group of the full $4$-shift. Up to reflection and the cellwise complement there are $88$ such automata, and we assign a status to each class. In $43$ classes the group contains every finite group; in $11$ it is $\langle\sh,F_f\rangle\cong\Z^2$; in $4$ it is $\langle\sh,F_f\rangle\times\Z/2$; in $2$ affine classes it is $\Z^2$ or $\Z^3$, strictly larger than $\langle\sh,F_f\rangle$; and $28$ classes are open. These entries rest on results of four companion papers: wall constructions, rigidity of the permutive rules and of the rules that read only two neighbours, edge moves for six classes, the rigidity theorem of Kitchens and Schmidt for the affine classes, and the groups of the classes of $94$ and $95$. We assemble them and add the following. An exhaustive SAT enumeration, repeated with several solvers, shows that in each open class every automorphism of radius at most $3$ lies in $\langle\sh,F_f\rangle$, and that for the rules $110$ and $30$ the block maps of radius at most $4$ that commute with $F_f$ are the maps $\sh^aF_f^b$ and two constants. The automaton $110$R has no wall of the known kinds, of any length. In at least $12$ classes, among them $10$ open ones, no configuration has a stuck edge, so the edge-move argument cannot start; computations, in part sampled, indicate where it stops in the other open classes.
\end{abstract}
